%% hypotheses_draft.tex
%% Standalone draft — not yet integrated into main paper.
%% Connects theory section predictions to testable hypotheses.

\section{Hypothesis Development}
\label{sec:hypotheses}

The theoretical framework in Section~\ref{sec:theory} generates three empirical
predictions linking political polarization to market reactions around auditor
change events. We formalize these as testable hypotheses below.

\subsection{Polarization and Market Disagreement}

When investors hold heterogeneous priors about auditor credibility, an identical
audit signal generates dispersed posterior beliefs. The cross-sectional dispersion
of posteriors determines both the magnitude of price movements and the volume of
trade \citep{kandel1995, harris1993}. Political polarization, by increasing the
dispersion of priors across ideological groups, amplifies these effects.

Our primary proxy for disagreement is \textit{absolute} cumulative abnormal
returns ($|CAR|$). Larger absolute returns around an event reflect a wider range
of investor valuations of the signal, consistent with greater belief
heterogeneity. Abnormal trading volume provides a complementary measure: under
heterogeneous beliefs, investors with sufficiently different valuations from the
market price find it optimal to trade \citep{harris1993}.

\begin{hyp}
    \label{h1}
    Investor reactions to auditor change disclosures are more extreme in
    politically polarized environments. Specifically, both the absolute
    cumulative abnormal return ($|CAR|$) and abnormal trading volume around
    auditor change events are increasing in political polarization.
\end{hyp}

Hypothesis~\ref{h1} tests whether the cross-sectional and time-series variation
in polarization — measured at the firm's headquarters state — predicts the
intensity of market reactions. A positive coefficient on the interaction of the
event indicator with the polarization measure, in specifications with firm and
year fixed effects, would be consistent with this prediction.

\subsection{Direction of the Price Reaction}

Hypothesis~\ref{h1} makes a prediction about the \textit{magnitude} of the
market reaction, not its direction. The direction depends on investors'
net assessment of the signal. However, our framework has an additional
implication for the \textit{signed} CAR. Polarization does not change the
average prior across investors; it increases the spread around the mean. If the
average prior is negative (investors interpret auditor changes as bad news on
average), polarization amplifies the negative reaction. If the average prior is
positive (e.g., for Big-4 upgrades), polarization amplifies the positive
reaction. We therefore expect:

\begin{hyp}
    \label{h2}
    Political polarization amplifies the signed CAR in the direction of
    the consensus view of the signal. For negative-signal events (e.g.,
    dismissals, downward quality changes), polarization increases the
    magnitude of negative abnormal returns. For positive-signal events
    (e.g., Big-4 upgrades), polarization increases the magnitude of
    positive abnormal returns.
\end{hyp}

Hypothesis~\ref{h2} requires splitting the sample by the likely direction of the
market signal, using the auditor change type (dismissal vs.\ resignation;
Big-4-to-non-Big-4 vs.\ non-Big-4-to-Big-4) as a proxy.

\subsection{Signal Ambiguity as a Moderator}

The model predicts that polarization effects are stronger when signals are more
ambiguous, i.e., when the likelihood function $P(s|\theta)$ is less precise.
Less precise signals leave more room for prior beliefs to determine posteriors,
amplifying the role of polarization. We proxy for signal ambiguity using
characteristics of the auditor change:

\begin{itemize}
    \item \textit{High ambiguity:} auditor changes with no disclosed reason,
          switches between auditors of similar size (non-Big-4 to non-Big-4 or
          Big-4 to Big-4), and changes that occur outside of fiscal year-end.
    \item \textit{Low ambiguity:} auditor changes accompanied by disclosed
          disagreements (Item 4.01 requires disclosure of any disagreements),
          Big-4 departures with no replacement, and changes coinciding with
          restatements or SEC investigations.
\end{itemize}

\begin{hyp}
    \label{h3}
    The effect of political polarization on absolute returns and trading volume
    is stronger for auditor changes that carry more ambiguous signals. The
    interaction of polarization with the event indicator is larger in magnitude
    for high-ambiguity events than for low-ambiguity events.
\end{hyp}

Hypothesis~\ref{h3} is tested by including a triple interaction of the event
indicator, polarization, and a signal-ambiguity indicator. A positive and
significant triple interaction coefficient supports the mechanism.
